% Generated draft: requires mathematical and source review.
\documentclass[11pt]{article}
\usepackage[margin=25mm]{geometry}
\usepackage{fontspec}
\setmainfont{Latin Modern Roman}
\usepackage{amsmath,amssymb,mathtools,unicode-math}
\setmathfont{Latin Modern Math}
\usepackage{enumitem,needspace,xurl,hyperref}
\setlength{\parindent}{0pt}
\setlength{\parskip}{5pt}
\setlength{\emergencystretch}{3em}
\newcommand{\N}{\mathbb{N}}
\newcommand{\Z}{\mathbb{Z}}
\newcommand{\Q}{\mathbb{Q}}
\newcommand{\R}{\mathbb{R}}
\newcommand{\C}{\mathbb{C}}
\newcommand{\F}{\mathbb{F}}
\newcommand{\A}{\mathbb{A}}
\newcommand{\PP}{\mathbb{P}}
\newcommand{\E}{\mathbb{E}}
\newcommand{\eps}{\varepsilon}
\newcommand{\dd}{\,\mathrm{d}}
\begin{document}
\section*{UnsolvedMath ID: 20000530}
\textbf{Source problem code:} catalogue entry
\subsection{Definitions and mathematical statement}
The supplied reconstruction uses the Euclidean unit ball \(B\subset\R^3\), the quadratic phase \(x_3^2\), and distinct lines \(V_j\) spanned by nonzero null vectors for
\(v_1^2+v_2^2-v_3^2=0\). For \(v_j\in V_j\setminus\{0\}\), define
\[
 I_\lambda(f_1,\ldots,f_N)=\int_B e^{i\lambda x_3^2}
 \prod_{j=1}^N f_j(x\cdot v_j)\,dx.
\]
Christ's intended question is whether, for arbitrary bounded measurable \(f_j:\R\to\C\), an estimate
\[
 |I_\lambda(f_1,\ldots,f_N)|\le C|\lambda|^{-\epsilon}
 \prod_{j=1}^N\|f_j\|_\infty\qquad(|\lambda|\ge1)
\]
holds for some \(\epsilon>0\), and what the optimal exponent is. The precise allowable dependence of \(C,\epsilon\) on \(N\) and the directions is not supplied. The extraction instead says ``different light cones'' and corrupts the phase, estimate symbol and epsilon; the verification stops after announcing the phase repair. Thus the displayed null-line/\(x_3^2\) formulation is the supplied research interpretation, not independently verified recovered text.

The reported partial theorem restricts the input class. For every fixed integer degree bound \(D\ge0\), every finite collection of nonzero null directions, and arbitrary real polynomials \(P_j\) of degree at most \(D\), it asserts
\[
 \left|\int_B e^{i\lambda x_3^2}\prod_{j=1}^N e^{iP_j(x\cdot v_j)}\,dx\right|
 \le C_D|\lambda|^{-1/2}\qquad(|\lambda|\ge1),
\]
uniformly in \(N\), placement of directions and all polynomial coefficients. Repeated directions are not excluded in this partial theorem. For finite complex measures \(\mu_j\) on the real degree-\(D\) polynomial coefficient space and mixtures
\(f_j(t)=\int e^{iP(t)}\,d\mu_j(P)\), the claimed bound becomes
\[
 |I_\lambda(f_1,\ldots,f_N)|\le C_D|\lambda|^{-1/2}
 \prod_j\|\mu_j\|_{\mathrm{TV}}.
\]
This is a product total-variation bound, not the original product \(L^\infty\)-norm estimate for arbitrary inputs. The report gives exponent \(1/2\) as sharp in this class because taking all inputs equal to one leaves a stationary-phase term of order \(|\lambda|^{-1/2}\).

The report also states an exact tensor criterion for quadratic cancellation. With \(e_3=(0,0,1)\), cancellation by quadratic ridge phases \(a_j(x\cdot v_j)^2\) is the matrix identity
\[
 \lambda e_3e_3^{\mathsf T}+\sum_j a_jv_jv_j^{\mathsf T}=0.
\]
This describes cancellation of the quadratic part, not a decay theorem for all bounded functions. The source's statement that estimate (0.2) had only been proved for \(N\le5\), with reference [4], is a historical qualification; no current verification or full bibliographic reference is supplied.
\subsection{Short English statement}
Seek oscillatory decay with arbitrary bounded ridge inputs along light-cone directions. The supplied reconstructed partial result proves sharp half-power decay for bounded-degree polynomial chirps and their finite-measure mixtures, a restricted class whose total-variation estimate does not resolve the original \(L^\infty\) question.
\subsection{Sources}
\begin{itemize}
\item Local UnsolvedMath registry, permanent ID \(20000530\), source code \texttt{catalogue entry}; AIM analysis corpus, \texttt{aim-analysis-notes.json}, notes[61]. Attribution: Ulam AI / UnsolvedMath Contributors, CC BY 4.0.
\item AIM \emph{Carleson theorems and multilinear operators}, Question 3, attributed to Michael Christ, equation (0.2): \url{https://aimath.org/pastworkshops/multilinopsproblems.pdf}. The source retains reference [4] for its \(N\le5\) claim; no bibliographic entry for [4] is supplied. Original literature and remarks fields are empty.
\item The verification says official PDF and TeX can repair the extraction but ends after ``the phase is''. The report reconstructs distinct null lines from ``different light cones''; this draft claims no new inspection.
\item Registry dated August 14, 2026: statement \texttt{reconstructed\_unverified}, status \texttt{partially\_solved}, classification \texttt{PARTIAL-PROGRESS}; attempt 1 records problem status open, partial result, rigor proved, novelty confidence low. The tensor/chirp/mixture contribution is a machine-generated research aid, not peer reviewed, requiring independent verification before citation.
\end{itemize}
\end{document}
